% TODO checklist:
% - [x] README.md: Platform overview, key features, CINECA context
% - [x] Q&A.md: Q1-Q10 (High-level understanding), Q50 (SOTA comparison)
% - [x] Abstract in sapthesis-doc.tex: Key motivations and gaps addressed

\section{Motivation and Context}
\label{sec:motivation}

The landscape of artificial intelligence has undergone a transformative shift with the advent of Large Language Models (LLMs). Models such as GPT-4~\cite{openai2023gpt4}, LLaMA~\cite{touvron2023llama}, and their successors have demonstrated unprecedented capabilities in natural language understanding, reasoning, and generation. These advances have catalyzed a new paradigm in software development: the construction of \emph{AI agents}---autonomous software entities capable of perceiving their environment, reasoning about goals, and executing multi-step actions to accomplish complex tasks~\cite{yao2022react}.

\subsection{The Rise of AI Agents in Enterprise Applications}

The progression from simple chatbots to sophisticated AI agents represents a fundamental evolution in how organizations leverage artificial intelligence. Unlike traditional conversational interfaces that merely respond to queries, AI agents can:

\begin{itemize}
    \item \textbf{Plan and decompose} complex tasks into manageable sub-tasks
    \item \textbf{Invoke external tools} to gather information or execute actions
    \item \textbf{Reason about intermediate results} and adapt their approach
    \item \textbf{Maintain context} across multi-turn interactions
    \item \textbf{Query structured knowledge} from databases and knowledge graphs
\end{itemize}

This capability has driven significant interest in deploying AI agents across enterprise domains, from customer service and document processing to scientific research and data analysis. Organizations increasingly seek to augment human expertise with AI systems that can autonomously handle routine inquiries, synthesize information from multiple sources, and provide actionable insights.

\subsection{The High-Performance Computing Context}

Research institutions and supercomputing centers face unique challenges in managing computational resources, scientific workflows, and the vast amounts of data generated by modern research. High-Performance Computing (HPC) environments, such as those operated by national supercomputing centers, support diverse user communities spanning bioinformatics, physics simulation, climate modeling, and many other domains.

In these environments, researchers and operators must navigate:

\begin{itemize}
    \item Complex job scheduling systems with intricate resource allocation policies
    \item Heterogeneous software stacks spanning multiple programming languages and frameworks
    \item Regulatory compliance requirements for data privacy and research integrity
    \item Multi-institutional collaborations requiring fine-grained access control
    \item The need to rapidly onboard new users and adapt to evolving research requirements
\end{itemize}

AI agents offer a compelling solution for many of these challenges by providing intuitive interfaces that abstract away underlying complexity while maintaining the precision and auditability that research environments demand.

\subsection{CINECA: Italy's National Supercomputing Center}

CINECA (Consorzio Interuniversitario del Nord Est per il Calcolo Automatico)~\cite{cineca2023} serves as Italy's premier supercomputing consortium, providing computational infrastructure and services to universities, research institutions, and industrial partners across the nation and beyond. Founded in 1969, CINECA has evolved to become one of Europe's leading HPC centers, operating some of the world's most powerful supercomputers and supporting thousands of researchers across diverse scientific disciplines.

The organization's mission encompasses:

\begin{itemize}
    \item Providing world-class computational resources for scientific research
    \item Developing and maintaining software tools for data management and analysis
    \item Supporting bioinformatics and life sciences research initiatives
    \item Enabling secure, multi-tenant access for diverse user communities
    \item Advancing the state of the art in high-performance computing applications
\end{itemize}

Within this context, there exists a compelling need for intelligent systems that can assist researchers in navigating complex computational workflows, querying interconnected research data, and managing resources efficiently. The intersection of AI agent capabilities with HPC requirements presents both significant opportunities and substantial technical challenges.

\subsection{The Gap Between Research and Production}

Despite the rapid proliferation of AI agent frameworks and tools, a significant gap persists between research-oriented prototypes and production-ready enterprise systems. Popular frameworks such as LangChain~\cite{langchain2023} and AutoGPT~\cite{significant2023autogpt} have demonstrated impressive capabilities in controlled settings, yet deploying these systems in enterprise environments reveals critical limitations:

\begin{enumerate}
    \item \textbf{Security and Governance}: Research frameworks often lack comprehensive authentication, authorization, and audit capabilities required for enterprise deployments.
    
    \item \textbf{Multi-Tenancy}: Most existing tools assume single-tenant deployments, making them unsuitable for organizations serving multiple user communities with distinct access requirements.
    
    \item \textbf{Reliability and Observability}: Production systems require robust monitoring, health checks, and graceful degradation---capabilities rarely prioritized in research-focused implementations.
    
    \item \textbf{Tool Integration Standards}: The absence of standardized protocols for tool integration leads to fragmented ecosystems and vendor lock-in.
    
    \item \textbf{Knowledge Graph Integration}: While many frameworks support document retrieval, few provide native integration with graph databases for structured knowledge querying.
\end{enumerate}

This thesis addresses these gaps by presenting the design and implementation of an enterprise-grade AI agent orchestration platform developed in collaboration with CINECA. The platform bridges the divide between academic agent research and industrial deployment requirements, providing a reference architecture for organizations seeking to deploy AI agents in production environments.

\subsection{The Model Context Protocol Opportunity}

The emergence of the Model Context Protocol (MCP)~\cite{anthropic2024mcp} presents a timely opportunity to address the tool integration challenge. MCP provides a standardized specification for how AI systems discover and invoke external tools, enabling interoperability across different agent implementations and reducing the fragmentation that has plagued the ecosystem.

By adopting MCP principles, the Cineca Agentic Platform establishes a foundation for extensible tool integration that can evolve with the rapidly advancing field of AI agents. This protocol-first approach ensures that investments in tool development remain valuable as the underlying AI technologies continue to improve.

\subsection{Chapter Summary}

This chapter has established the motivation for developing an enterprise-grade AI agent orchestration platform by:

\begin{itemize}
    \item Describing the transformative potential of LLM-powered AI agents
    \item Contextualizing the unique requirements of HPC and research environments
    \item Introducing CINECA as the institutional setting for this work
    \item Identifying the critical gaps between research frameworks and production requirements
    \item Highlighting the opportunity presented by emerging tool integration standards
\end{itemize}

The following sections of this introductory chapter will elaborate on specific use cases, stakeholder requirements, the formal problem statement, research objectives, and the contributions of this thesis.

